\documentclass[ecta]{econsocart}
\usepackage{amsmath,amssymb,amsthm,booktabs,tabularx,natbib,microtype}
\usepackage{hyperref}
\hypersetup{colorlinks=true,allcolors=blue}
\emergencystretch=2em
\startlocaldefs
\newtheorem{theorem}{Theorem}
\newtheorem{proposition}{Proposition}
\newtheorem{assumption}{Assumption}
\newtheorem{lemma}{Lemma}
\theoremstyle{remark}
\newtheorem{remark}{Remark}
\newtheorem{example}{Example}
\newcommand{\R}{\mathbb R}
\newcommand{\E}{\mathbb E}
\endlocaldefs
\begin{document}
\begin{frontmatter}
\title{Neural Bellman Operators}
\runtitle{Neural Bellman Operators}
\begin{aug}
\author[id=au1,addressref={add1}]{\fnms{Qian}~\snm{QI}\ead[label=e1]{qiqian@pku.edu.cn}}
\address[id=add1]{\orgname{Peking University}}
\end{aug}
\begin{abstract}
This paper develops Neural Bellman Operators for stochastic economic models with endogenous preferences, recursive utility, and strategic interaction. Policy evaluation and feasible improvement are separated, and the enacted policy is evaluated independently of its critic. We establish boundary-aware continuous-time and finite-horizon bounds that convert evaluation defects and feasible improvement gaps into economic value loss. For a preference--portfolio economy with an explicit liquidation contract, we prove that a higher adjustment cost reduces the optimal discounted preference-adjustment budget. Executed neural updates recover the continuous-time Merton benchmark and approximate a nonlinear stochastic preference economy; three completed neural runs have maximum date-zero policy value losses between .0113 and .0236 on the declared finite model. A re-solved cost panel exhibits an active preference margin and the predicted adjustment-budget ordering. Recursive-utility, temporal-self, dynamic-game, and correlated dynamic-control calculations distinguish the mathematical and computational requirements of the extensions. Independent error definitions expose the remaining diffusion-discretization and portfolio-policy discrepancies instead of attributing them to a small training residual.
\end{abstract}
\begin{keyword}
\kwd{Stochastic control}\kwd{Neural policy iteration}\kwd{Endogenous preferences}\kwd{Recursive utility}\kwd{Policy verification}
\end{keyword}
\end{frontmatter}
\section{Introduction}
Dynamic economic models increasingly place the formation of preferences, the allocation of financial risk, and the adjustment of physical resources inside the same decision problem. The relevant computational object is then a controlled operator: a policy changes both the payoff and the law under which continuation utility is evaluated. A useful approximation must preserve this distinction. In particular, reducing a joint scalar training objective need not improve an economic policy, and a small error in a value representation need not identify the policy at a nearly flat optimum.

This paper develops Neural Bellman Operators (NBO) as a separated evaluation--improvement method for such problems. The critic approximates the continuation value of a specified feasible policy. The actor improves that policy against a frozen continuation object. A third, independent calculation evaluates the enacted policy and its feasible deviation gains. We study both the continuous generator formulation and a positive finite-horizon operator formulation. The latter makes the stochastic transition and the action maximization explicit and provides a reference against which a trained neural representation can be assessed.

There are three substantive results. First, we give a quantitative connection between policy evaluation, boundary approximation, feasible improvement, and loss of economic value. The continuous-time bound applies to a stopped controlled diffusion and requires uniform, rather than merely sampled, residual bounds. Its finite-horizon counterpart is directly computable on a finite state--action approximation. Second, we establish a comparative-static result for endogenous preference formation. When the adjustment cost enters as a price of a discounted adjustment budget, an increase in that price reduces the optimal budget. This conclusion holds for any optimizers and does not require differentiability of the value function or a pointwise sign restriction on preference adjustment. Third, we implement the separated method, including automatic differentiation of a continuous HJB benchmark and trained nonlinear actor and critic networks for a stochastic preference--portfolio problem. Independent policy evaluation separates value accuracy from agreement of the full policy vector.

The economic application preserves the interior preference and wealth dynamics of the original model. It completes the model at the boundary by an explicit early-liquidation contract. The preference coefficient has nonzero diffusion, so a bounded drift cannot keep it in a closed interval with probability one. Likewise, the original consumption floor is incompatible with a self-financing wealth process constrained to the lower endpoint. Liquidation is an economically different boundary condition from reflection or a capital injection; it is specified as such, including the settlement fee and its sensitivity. No transition replenishes wealth and then continues consumption.

The numerical evidence has distinct roles. A continuous-time Merton benchmark tests the HJB derivatives and the separated gradient graph. A two-state NDU problem tests a trained nonlinear neural representation against an independent exhaustive stochastic Bellman solver. A re-solved adjustment-cost panel evaluates the economic comparative static. Recursive utility, sophisticated temporal selves, and a dynamic capacity game exercise additional operator interfaces. A coupled stochastic linear--quadratic problem checks dynamic cross-state dependence at dimensions 4, 8, and 16. These experiments retain the general economic scope without conflating an algebraic identity, a finite-model equilibrium, and a trained solution of a continuous nonlinear economy.

\subsection{Relation to the literature}
The closest methodological antecedents already combine policy improvement and neural approximation. \citet{kim2025} study physics-informed neural policy iteration with policy-error analysis. \citet{lee2024} combine Hamilton--Jacobi policy iteration with deep operator learning. \citet{cai2025} learn value and control through a martingale formulation without requiring an explicit Hamiltonian optimizer. Direct neural approximation of stochastic control policies is developed by \citet{han2016}. Continuous-time policy improvement and its hypotheses are analyzed by \citet{jacka2015}. Accordingly, neural approximation, operator factorization, and the absence of an analytical maximizer are not claimed here as inventions in themselves.

The contribution relative to these approaches is the particular combination of boundary-aware value-loss control, an economically interpretable preference-adjustment result, and an executable separation of the evaluation, improvement, and policy-certification objects. The finite-operator realization also permits a distinction between errors of neural representation and errors of the numerical economy. The bounds below are not presented as a priority claim for residual-based verification in general; they specify the conditions and constants used by this implementation and its application.

Our economic formulation connects portfolio choice \citep{merton1969}, recursive preferences \citep{epstein1989,duffie1992}, and the motivating endogenous-preference model of \citet{qi2025}. Grid and projection methods remain important comparators in computational economics \citep{judd1998,brumm2017}. In our two-state experiment, exhaustive backward induction is faster than neural training. That result is reported, not converted into an unobserved scalability advantage. A neural representation can nevertheless be evaluated as a function approximation to a controlled operator, and an independently measured value loss supplies a meaningful test of that representation.

The distinction from Deep BSDE methods \citep{han2018} is about the equations and training objects, not about whether an analyst can write a closed-form control. Semilinearity concerns dependence on the highest-order derivative. Eliminating a control analytically can leave nonlinear Hessian dependence. Conversely, an admissible BSDE representation need not be a Pontryagin adjoint system. Section~\ref{sec:trace} states the differentiation costs conditionally, including the diffusion rank and the accuracy demanded of a trace approximation.

\section{The economic problem and its boundary}
\label{sec:problem}
Let $S$ be a Markov state in an open bounded domain $D\subset\mathbb R^d$, let $A(s)$ be a compact feasible action correspondence, and let $\pi$ be an admissible feedback. Up to $\tau=T\wedge\inf\{r\geq t:S_r\notin D\}$, the state satisfies
\begin{equation}
dS_r=b(r,S_r,\pi_r)\,dr+\Sigma(r,S_r,\pi_r)\,dW_r.
\label{eq:state}
\end{equation}
Admissibility includes existence of the stopped process and integrability of the stated payoff. For a time-separable payoff with nonnegative discount rate $\rho$, define
\begin{equation}
J^\pi(t,s)=\mathbb E_{t,s}\left[\int_t^\tau e^{-\rho(r-t)}\ell(r,S_r,\pi_r)\,dr
 +e^{-\rho(\tau-t)}g(\tau,S_\tau)\right],\qquad V^*=\sup_\pi J^\pi.
\label{eq:payoff}
\end{equation}
The settlement $g$ is data of the problem, not an output of a projection routine. At $T$, it equals the terminal endowment utility. The operator is
\[
\mathcal L^a v=b^a\cdot\nabla v+\tfrac12\operatorname{tr}(\Sigma^a(\Sigma^a)'D^2v),
\quad R_a(v)=v_t+\mathcal L^a v+\ell^a-\rho v.
\]
Where a classical representation is valid, fixed-policy evaluation has $R_\pi(V^\pi)=0$, while optimality has $\sup_a R_a(V^*)=0$. Both have the same declared parabolic-boundary data. A reflected model would instead include regulators and an appropriate derivative boundary condition; the results are not obtained by replacing that condition with a value penalty.

\begin{assumption}[Well-defined evaluation and verification]\label{ass:verify}
The feasible controls yield a stopped state process; the payoff and stochastic integrals used below are integrable. For each fixed policy under discussion, evaluation is well defined. Comparison and the relevant verification inequalities hold for the specified generator and boundary operator. Whenever derivatives are used, the test function belongs to $C^{1,2}$ up to the relevant stopped boundary and its discounted It\^o formula is valid.
\end{assumption}
These hypotheses state the mathematical object to which a residual refers. They do not assert that an arbitrary smooth network or an arbitrary compact computational rectangle verifies them.

\subsection{Recursive utility and Markov closure}
For recursive preferences, fixed-policy utility is defined before optimization. On a domain of its value argument where the driver is admissible, it solves
\begin{equation}
-dY_r^\pi=f(r,S_r,\pi_r,Y_r^\pi)\,dr-Z_r^\pi dW_r,
\qquad Y_\tau^\pi=g(\tau,S_\tau).
\label{eq:bsde}
\end{equation}
Existence, uniqueness, comparison, and integrability of this policy-wise equation are prerequisites for taking its supremum. The Markov state must contain any variable on which the driver or future opportunities depend. A general continuous martingale is not replaced by a Gaussian shock merely because its quadratic variation is known; the executed models here use the stated Brownian or finite Markov specifications. The recursive extension and its sign domain are developed in Section~\ref{sec:recursive}.

\section{Separated Neural Bellman Operators}
\label{sec:method}
In the continuous realization, a critic $v_\xi$ is trained for a fixed enacted policy by its evaluation residual and the appropriate boundary loss. The actor parameters $\psi$ are updated against a detached critic jet:
\begin{align}
L_E(\xi;\pi)&=\mathbb E_\nu R_\pi(v_\xi)^2+\lambda_B\mathbb E_{\nu_B}|v_\xi-g|^2,\label{eq:critic}\\
L_I(\psi;\xi)&=-\mathbb E_{\nu_I}\mathcal H^{\pi_\psi}[\operatorname{stopgrad}(v_\xi,\nabla v_\xi,D^2v_\xi)].\label{eq:actor}
\end{align}
Feasibility is imposed by the actor parameterization or by explicit constrained maximization. A penalty for a general constraint is an approximation and must be accompanied by a feasibility error. During actor improvement no parameter gradient reaches $\xi$. In a game no player-$i$ gradient reaches the parameters of a frozen rival.

The actor loss is a local improvement surrogate. It is not generally the lifetime-objective gradient under an arbitrary sampling measure. Such an identification requires the appropriate occupancy or adjoint weighting and differentiability conditions. Likewise, minimizing a weighted sum of \eqref{eq:critic} and \eqref{eq:actor} over both parameter blocks solves a different optimization problem. Appendix~\ref{app:limits} retains the composite-loss and viscosity counterexamples and states what the approximation argument does establish.

\subsection{A positive finite-operator realization}
For the numerical economy, let $\mathcal S_h$ be a finite state set and $A_h$ a finite feasible action set. Liquidation proceeds are included in $r_n^a$ and surviving transitions in a substochastic kernel $P_n^a$. With $\delta=e^{-\rho\Delta}$,
\begin{equation}
T_n^a w=r_n^a+\delta P_n^a w,\qquad T_n w=\max_{a\in A_h}T_n^a w.
\label{eq:finite}
\end{equation}
The factor $\delta$ applies to surviving full steps; an intra-step liquidation has its own fractional-step discount in $r_n^a$. Boundary states are settled immediately.

The finite NBO algorithm works backward. It freezes the next critic, computes each feasible action's continuation using the declared stochastic operator, and trains the current actor against these action values. It then trains the current critic to evaluate the actor's enacted hard action. The numerical experiment uses a cross-entropy term for the maximizing action of the \emph{neural continuation}, together with its softmax-weighted action-value regret. Optimal reference values or labels never enter this training loss. At deployment the actor chooses its hard argmax, not a mixture. All continuation targets are detached. This is a finite-horizon neural Bellman factorization, not a claim that one Adam update performs exact policy iteration.

The reference solver is an independently coded NumPy backward induction with exhaustive action evaluation and positive bilinear interpolation. The neural operator is implemented separately in PyTorch. An operator-agreement test precedes training. The trained policy is then re-evaluated by the reference operator; this evaluation is different from reading the critic's prediction. All maxima used for finite-model certificates range over the entire stated finite set.

\subsection{Exact improvement and quantitative approximation}
\begin{proposition}[Exact separated-operator limit]\label{prop:exact}
Let the policy space $\Pi$ be compact in the uniform topology. Let evaluation $E$ be continuous into a precompact subset of $C^{1,2}$, endowed with the uniform norm on the function, its time derivative, gradient, and Hessian. Assume a maximizing selector $I$ maps $\Pi$ into itself, is continuous at accumulation points, and satisfies the comparison and verification conditions of Assumption~\ref{ass:verify}. If $\pi_{j+1}=I(\pi_j)$, every accumulation point $p$ satisfies $E(p)=V^*$, and the common limiting value solves the HJB boundary problem.
\end{proposition}
The decisive step is equality of limiting \emph{values}, not equality of successive limiting policies. The full proof is in Appendix~\ref{app:proofs}. Closure under the exact selector is a real hypothesis and is not automatic for a finite neural class.

\begin{theorem}[Boundary-aware policy-value certificate]\label{thm:continuous}
Under Assumption~\ref{ass:verify}, let $v$ be a test function and $\pi$ a feasible policy. Suppose, uniformly on the stopped domain,
\[
|R_\pi(v)|\leq\epsilon,\qquad
0\leq\sup_{a\in A(s)}\{R_a(v)-R_\pi(v)\}\leq\eta,
\qquad |v-g|\leq b\quad\hbox{on the parabolic boundary}.
\]
Then, with $C_\rho(h)=(1-e^{-\rho h})/\rho$ and $C_0(h)=h$,
\begin{equation}
0\leq V^*(t,s)-J^\pi(t,s)\leq 2b+C_\rho(T-t)(2\epsilon+\eta).
\label{eq:continuous-bound}
\end{equation}
No smoothness of $V^*$, and no claim that an actor stationary point is a maximizer, is required for this inequality.
\end{theorem}
The boundary term is not discounted at the terminal date: exit can occur earlier. Theorem~\ref{thm:continuous} turns three separately meaningful errors into a value loss. A sampled mean-square loss does not supply its uniform hypotheses. A verified residual Lipschitz constant and a covering radius can bridge a finite audit to a uniform bound; such constants must be established, not inferred from a plot.

\begin{theorem}[Finite-horizon neural-policy certificate]\label{thm:finite}
For the finite operators \eqref{eq:finite}, let $\widehat v_N$ have terminal error at most $b$. For an enacted policy $\pi$, define
\[
e_n=\|\widehat v_n-T_n^\pi\widehat v_{n+1}\|_\infty,
\qquad \eta_n=\|T_n\widehat v_{n+1}-T_n^\pi\widehat v_{n+1}\|_\infty.
\]
Then
\begin{equation}
\|V_n^*-V_n^\pi\|_\infty
\leq 2\delta^{N-n}b+\sum_{j=n}^{N-1}\delta^{j-n}(2e_j+\eta_j).
\label{eq:finite-bound}
\end{equation}
If the enacted policy is re-evaluated exactly on this finite model, take $\widehat v=V^\pi$; then $e_j=b=0$ and the same bound uses only the independently evaluated feasible improvement gaps.
\end{theorem}
This theorem bounds a finite model, not automatically the underlying diffusion. Time, interpolation, quadrature, action coverage, and boundary-crossing errors remain distinct. An additional verified comparison between the finite and continuous economies would be needed to add them to \eqref{eq:finite-bound}.

\subsection{Global action improvement versus parameter stationarity}
The Hamiltonian $H(a)=-(a^2-1)^2+.2a$ on $[-2,2]$ has a stable local ascent attractor near $-.973994$ and a global maximum near $1.024120$. Their payoff gap is $.3998746$. Both stationary points are isolated. Thus an exact fast critic, diminishing two-timescale steps, and isolated actor stationary points do not imply a policy-improvement fixed point.

A usable sufficient condition concerns the \emph{action} maximization. If $H$ is differentiable and concave on a convex compact feasible set, then
\begin{equation}
H^*-H(a_0)\leq\max_{a\in A}\langle\nabla H(a_0),a-a_0\rangle.
\label{eq:fwgap}
\end{equation}
The right-hand side is a global supporting-hyperplane gap and handles boundary optima. Pointwise concavity does not imply concavity in shared neural parameters. Our finite implementation therefore audits the maximizing finite action directly. The separate asymptotic stochastic-approximation statement, with its invariant-set conclusion and additional optimality conditions, is given in Appendix~\ref{app:limits}.

\section{Endogenous preferences and portfolio choice}
\label{sec:ndu}
The preference state $u$ is the coefficient of relative risk aversion and $X$ is financial wealth. The controls are consumption $c$, preference adjustment $\theta$, and the risky portfolio share $\pi$. Before liquidation,
\begin{align}
du&=\theta\,dt+.05\,dZ^u,\\
dX&=[(.02+.06\pi)X-c]dt+.2\pi X\,dZ^X,
\qquad d\langle Z^u,Z^X\rangle=-.25\,dt.\label{eq:ndu-state}
\end{align}
The domain is $(1.2,2.8)\times(.5,2)$ and the control box is
$[.05,.8]\times[-.2,.2]\times[-.5,.8]$. The horizon is one year and $\rho=.04$. Utility is cardinally fixed as
\[
\ell_k(u,c,\theta)=\frac{c^{1-u}}{1-u}-\frac{k}{2}\theta^2,
\qquad G(u,X)=-.02(u-2)^2+.1\log X.
\]
Consumption units and this normalization are held fixed throughout the cost panel. An affine transformation that depends on the controlled preference state generally changes the economic problem.

\subsection{The liquidation contract}
The process stops when it first reaches the domain boundary or the horizon. Settlement is
\begin{equation}
g(t,u,X)=G(u,X)-F(1-t),\qquad F=8.
\label{eq:settlement}
\end{equation}
The early termination fee represents the utility cost of closing the investment and preference-adjustment arrangement before maturity. This is an explicit covenant/liquidation specification, not a calibrated estimate of such a fee. Its size is varied separately below. The fee vanishes at maturity and does not depend on $k$. No additional wealth is supplied after settlement because the problem has terminated.

This completes rather than disguises the boundary economics. Constant preference diffusion rules out the earlier viable-rectangle interpretation. At $X=.5$, even the largest wealth drift under the original action box is $-.016$. Keeping the original interior process and the original control floor therefore requires exit, reflection with an explicit regulator, or a different resource model. We choose exit. Economic statements about this contract must not be interpreted as statements about a silently reflected self-financing economy.

The interior HJB is
\begin{align}
0=V_t+\sup_{c,\theta,\pi}\bigg\{&\frac{c^{1-u}}{1-u}-\frac{k}{2}\theta^2
+\theta V_u+[(.02+.06\pi)X-c]V_X-\rho V\nonumber\\
&+\tfrac12(.05)^2V_{uu}+\tfrac12(.2\pi X)^2V_{XX}
-.25(.05)(.2)\pi X V_{uX}\bigg\},\label{eq:ndu-hjb}
\end{align}
with \eqref{eq:settlement} on the parabolic boundary. All three second-order contributions are retained.

\subsection{Policies and economic interpretation}
Where the maximizer is interior and the indicated derivatives have the required signs,
\begin{equation}
c=V_X^{-1/u},\qquad \theta=V_u/k,\qquad
\pi=-\frac{.06V_X}{.2^2XV_{XX}}-
\frac{.05(-.25)V_{uX}}{.2XV_{XX}}.
\label{eq:foc}
\end{equation}
These are candidates, not substitutes for constrained maximization. Consumption and adjustment have concave action objectives. The portfolio objective is concave only when $V_{XX}\leq0$; otherwise endpoint comparisons are essential. Appendix~\ref{app:ndu} gives the full selection rule, including zero curvature and nonpositive marginal wealth value.

The cross derivative produces a hedge against preference shocks correlated with returns. Its sign is endogenous: neither a negative shock correlation nor an assumed positive $V_{uX}$ alone establishes a global portfolio comparative static. Likewise, dividing $V_u$ by $k$ does not determine how the re-solved derivative changes with $k$. The following result gives a global economic conclusion without either shortcut.

\begin{theorem}[The price of preference adjustment]\label{thm:cost}
Suppose the feasible policy set, stopping rule, and settlement contract do not depend on $k$. Write
\[
A(\pi)=\mathbb E\left[\int_t^\tau e^{-\rho(r-t)}\frac{c_r^{1-u_r}}{1-u_r}\,dr
+e^{-\rho(\tau-t)}g(\tau,S_\tau)\right],
\quad B(\pi)=\mathbb E\int_t^\tau e^{-\rho(r-t)}\frac{\theta_r^2}{2}\,dr.
\]
Then $V(k)=\sup_\pi\{A(\pi)-kB(\pi)\}$ is convex and nonincreasing. If $k_1<k_2$ and optimal policies $\pi_1,\pi_2$ exist, $B(\pi_2)\leq B(\pi_1)$. At a differentiability point with an attained optimizer, $V'(k)=-B(\pi_k)$. The flexible-preference value weakly exceeds the value under the feasible restriction $\theta\equiv0$.
\end{theorem}
The theorem concerns the optimal discounted adjustment budget, not a pointwise ordering of $\theta(t,u,X)$. The same proof applies to the finite stochastic economy, so the numerical panel tests a precise economic restriction rather than an unverified policy-surface narrative.

\section{Numerical implementation and results}
\label{sec:results}
\subsection{Continuous HJB benchmark}
We first solve the stationary Merton economy with CRRA coefficient 2, discount $.04$, riskless return $.02$, excess return $.06$, and volatility $.2$. A homothetic critic is $v(x)=-\exp(w)/x$ and a two-logit actor represents $m=c/x\in[.005,.07]$ and $\pi\in[-.5,1]$. First and second wealth derivatives are obtained by automatic differentiation. Critic optimization minimizes the scaled fixed-policy HJB residual. Actor optimization holds these derivatives fixed. Neither loss receives the analytical optimal controls or value as a label.

Independent evaluation gives $A^\pi=1/(mD)$ for $J^\pi=-A^\pi/x$, where
$D=.04+.02+.06\pi-m-.04\pi^2>0$. The optimum has $m^*=.04125$, $\pi^*=.75$. Across seeds 0, 1, and 2 the executed actor recovers $m$ within $8.25\times10^{-8}$ and $\pi$ within $8.29\times10^{-10}$. The maximum independently evaluated value loss on $[.5,2]$ is $4.70\times10^{-9}$. The scaled HJB residual is at most $1.38\times10^{-8}$. A graph assertion verifies that actor backpropagation leaves the critic gradient disconnected. This is a homothetic neural benchmark with three trained parameters, not evidence for a deep high-dimensional approximation theorem. Its value is that the continuous generator, automatic derivatives, evaluation, and actor graph can all be checked independently.

\subsection{Stochastic NDU transition and neural architecture}
For the preference--portfolio model, four equiprobable branches use independent signs $z_1,z_2\in\{-1,1\}$:
\begin{align}
\Delta u&=\theta\Delta+.05\sqrt\Delta z_1,\\
\Delta X&=[(.02+.06\pi)X-c]\Delta
+.2\pi X\sqrt\Delta\bigl[-.25z_1+\sqrt{1-.25^2}z_2\bigr].
\label{eq:quadrature}
\end{align}
If the straight segment exits, it settles at its first intersection and terminates; otherwise continuation uses positive bilinear interpolation. The running payoff is integrated with the exact discount factor for the elapsed fractional step while holding the current flow fixed. This defines a killed-Euler numerical economy, not exact Brownian first-passage integration. Interior drift and covariance tests are run before boundary corrections.

The nonlinear actor and critic each have two hidden layers of width 48 with tanh activations. The actor has 125 outputs, representing five values for each of the three controls. The critic uses normalized states and log distances to the four boundary faces; it is represented as $G+(1-t)N_\xi$ away from boundary nodes, with exact prescribed settlement at boundary nodes. The finite model does not require this boundary-switched representation to be a globally $C^2$ function. Consequently, the numerical certificate applied to it is Theorem~\ref{thm:finite}, not an unverified use of Theorem~\ref{thm:continuous}.

At each of six dates, training uses all interior nodes of a $25\times31$ lattice, 500 actor Adam steps and 1,000 critic Adam steps, followed by critic L-BFGS refinement. The learning rate is $.004$. Next-date critics and actor targets are frozen. This experiment is a lattice-trained finite-operator realization, not a mesh-free sampling experiment. Seeds 0, 1, and 2 are reported without selecting the best seed.

\begin{table}[ht]
\centering
\caption{Independent evaluation of the trained NDU policy}
\label{tab:neural}
\begin{tabular}{lrrr}
\toprule
Quantity & Seed 0 & Seed 1 & Seed 2\\
\midrule
Maximum date-zero value loss & .011291 & .023526 & .012797\\
Mean interior date-zero value loss & .001583 & .001803 & .001540\\
Maximum critic--policy value difference & .144195 & .070268 & .104457\\
Finite gain-based upper bound & .259331 & .264491 & .520503\\
Training seconds & 19.039 & 19.609 & 20.737\\
\bottomrule
\end{tabular}
\par\smallskip\footnotesize Values compare the independently re-evaluated hard actor with exhaustive optimization on the identical finite stochastic economy. The gain bound uses maxima at all dates and is conservative. Times are local single-thread CPU measurements, not universal complexity constants.
\end{table}

All three completed runs meet the stated finite-model date-zero value-loss target of $.05$. The critic's own error is larger than the enacted policy loss, illustrating why a critic prediction is not a policy-evaluation certificate. Boundary discrepancies are at most $4.80\times10^{-7}$ in the float32 neural output. For seed 0 the maximum component differences in $(c,\theta,\pi)$ from the finite optimal action are $(.375,.2,.975)$; the corresponding interior mean differences are $(.00183,.00335,.11028)$. The portfolio policy is therefore not uniformly accurate even when value regret is small. Every seed's complete component vector is retained in the ledger.

Exhaustive reference backward induction takes about $.60$ seconds in the seed-0 comparison, versus $19.04$ seconds for training and $.65$ seconds for independent policy evaluation. There is no measured speed advantage in this two-state task. The result instead establishes that a genuinely trained actor and critic implement the stated stochastic operator with small independently evaluated finite-model policy loss. Whether representation reuse or higher-dimensional nonlinear models reverse the cost comparison requires matched experiments on those objects.

The initial neural pilot failed: its maximum policy value loss was $4.10$, its critic error $4.49$, and its boundary representation was inaccurate. The failure and its executed source are retained. Two attempted twelve-date runs reached training output but timed out before completing their final audit; they are recorded as incomplete, not successful runs. These records are not included in the three-seed success count.

\subsection{Re-solved preference-adjustment costs}
The cost panel uses twelve dates, the $25\times31$ state lattice, and the 125-action stochastic operator. Each $k$ is independently re-solved, along with the restricted $\theta=0$ economy. Table~\ref{tab:cost} evaluates $(u,X)=(2,1.25)$.
\begin{table}[ht]
\centering
\caption{Costly preference formation under the liquidation contract}
\label{tab:cost}
\begin{tabular}{lrrr}
\toprule
 & $k=.5$ & $k=2$ & $k=8$\\
\midrule
Value & $-1.351522$ & $-1.372708$ & $-1.404420$\\
Discounted adjustment budget $B$ & .016520 & .011921 & .002309\\
Value of permitting adjustment & .064480 & .043293 & .011582\\
Date-zero adjustment $\theta$ & .2 & .2 & .1\\
Exit probability & .248935 & .249938 & .249920\\
Interior share with nonzero $\theta$ & .896927 & .764993 & .475637\\
\bottomrule
\end{tabular}
\par\smallskip\footnotesize The budget includes the factor $\theta^2/2$ and time discounting. Exit probabilities refer to this numerical liquidation contract. Nonzero-action shares average over all decision dates and interior lattice states, not over a population.
\end{table}
The budget falls by approximately 86 percent between $k=.5$ and $k=8$. Value monotonicity, budget monotonicity, and dominance over the zero-adjustment policy hold at all audited finite states and dates, up to the implementation tolerance. The preference margin is economically active rather than grid-locked. The result is conditional on the cardinal normalization and settlement contract; it is not identification of these primitives from observed consumption or portfolios.

\subsection{Discretization and boundary sensitivity}
Separate refinements are deliberately not labeled solution-error certificates. At the same central state, values for $(N,n_u,n_X,n_a)$ are
\[
\begin{array}{c|r}
(6,13,16,5)&-1.423383\\
(12,13,16,5)&-1.564151\\
(12,25,31,5)&-1.372708\\
(24,25,31,5)&-1.451486\\
(24,49,61,5)&-1.344458\\
(12,25,31,7)&-1.359628
\end{array}
\]
Here $n_a$ is the number of values per control. The changes are material. Refining time alone while holding a positive interpolation grid fixed need not improve a diffusion approximation: the interpolation contribution to a generator can scale as $h^2/\Delta$. Action refinement and boundary-crossing approximation also matter. We therefore report a verified finite-model neural approximation and a re-solved finite-model cost panel, but not a numerically certified continuous-diffusion value error. The analytical comparison in Theorem~\ref{thm:cost} does apply to the continuous contract under its stated existence conditions.

Changing the settlement fee from 8 to 6 or 10, with the twelve-date baseline discretization held fixed, changes central value to $-1.363909$ or $-1.379622$. These are different contracts rather than alternative numerical corrections to the same contract. The complete exit and adjustment calculations are retained so that subsequent boundary refinements can be evaluated without treating post-correction feasibility as a zero-error result.

\section{Recursive utility and equilibrium extensions}
\label{sec:extensions}
\subsection{Recursive utility}\label{sec:recursive}
Let $a=1-1/\psi$, $b=1-\gamma$, and $bV>0$. In the normalization used here, the Epstein--Zin continuous-time aggregator is
\begin{equation}
f(c,V)=\frac{\rho}{a}c^a(bV)^{1-a/b}-\frac{\rho b}{a}V,
\label{eq:ez}
\end{equation}
with the logarithmic limit when $a=0$. The architecture $V=\exp(W)/b$ enforces the sign domain when $b\ne0$. A compact positive margin for $bV$, along with consumption bounds, supplies local driver regularity; a sign assertion without executing the transform is not a domain audit.

The full portfolio evaluation equation is $V_t+\mathcal L^{c,\pi}V+f(c^\pi,v)=0$. Formal stationary homothetic ratios at $(\rho,\gamma,\psi,\mu,r,\sigma)=(.04,5,1.5,.08,.02,.2)$ are $\pi=.3$ and $m=.0455$. An algebraic ratio is distinct from a verified infinite-horizon solution. We retain these anchors and add an executed finite-horizon implicit recursive-utility calculation. The latter has horizon $.5$, ten dates, 101 wealth points and 80 consumption fractions, and solves each fixed-action implicit value equation on its negative-value branch. Its maximum implicit-equation residual is $2.46\times10^{-13}$; the minimum $bV$ is $.0625$; the executed terminal sign transform differs from the terminal target by at most $8.89\times10^{-16}$. Its continuation outside the wealth grid uses the specified homothetic numerical boundary data. This calculation is a recursive reference, not a trained neural portfolio solution. Appendix~\ref{app:recursive} supplies the equation, normalization, and verification requirements for the broader recursive application.

\subsection{Sophisticated temporal selves}
For discrete dates and $\delta=e^{-\rho\Delta}$, a sophisticated acting self takes the continuation policy of future selves as given:
\begin{align}
c_n^*(s)&\in\arg\max_c\{u(c)\Delta+\beta\delta\mathbb E[V_{n+1}(F(s,c,\varepsilon))]\},\\
V_n(s)&=u(c_n^*(s))\Delta+\delta\mathbb E[V_{n+1}(F(s,c_n^*(s),\varepsilon))].
\label{eq:selves}
\end{align}
Only improvement carries the extra factor $\beta$. For two consumption dates followed by terminal log wealth, zero interest and $\delta=1$, ordinary continuation has coefficient 2 on log wealth at the first decision. The sophisticated initial consumption share is $1/(1+2\beta)$. At $\beta=.7$ this is $.4166667$, whereas geometric-$\beta$ evaluation gives $.4566210$. The corrected executed evaluation residual is below $8.89\times10^{-16}$ and a dense feasible one-shot search finds no positive gain. Strict concavity and the analytical first-order condition establish global optimality for this particular interior log example; the finite search alone would provide only a lower bound on exploitability. The $\beta=1$ check is retained but is not used to detect an error that vanishes at $\beta=1$.

\subsection{Dynamic capacity competition}
For player $i$, a frozen-rival policy defines both the evaluation operator and an independent unilateral best-response problem. The static Cournot benchmark has Nash output $1/3$ and joint-profit output $1/4$; an automatic-differentiation test confirms that the rival gradient is absent from player $i$'s actor update.

We also solve a genuinely dynamic finite game. Two firms have capacities in $\{0,.5,1\}$, output in half-unit increments up to capacity, and investment $i\in\{0,1\}$ that adds $.5$ capacity when feasible. Demand intercept is $z\in\{2.5,3.5\}$ with transition matrix having diagonal $.8$ and off-diagonal $.2$. Per-period profit is $q_i(z-q_i-q_j)-.2i$, discount is $.96$, and the horizon is three production dates. Backward enumeration selects a pure Nash equilibrium in every state; lexicographic selection resolves multiplicity in three state-date problems. The maximum feasible one-shot deviation gain is computed as zero in this finite game. The full profile and selection rule are stored, including asymmetric equilibria. This is not a neural dynamic-game training result or a competitive-limit calculation. The latter would require a market-size normalization and a separately solved limiting economy; the operator formulation and its economic questions are retained in the appendix.

\section{Stochastic traces and coupled dynamics}
\label{sec:trace}
For symmetric $A$ and Gaussian probes $z_k$, write $q_k=z_k'Az_k$ and $\bar q_K=K^{-1}\sum q_k$. The squared-residual objective satisfies
\begin{equation}
\mathbb E(a-\bar q_K/2)^2=(a-\operatorname{tr}A/2)^2+\frac{\|A\|_F^2}{2K}.
\label{eq:trace}
\end{equation}
Averaging $(a-q_k/2)^2$ instead leaves a variance penalty independent of $K$. With $A=\left(\begin{smallmatrix}2&1\\1&3\end{smallmatrix}\right)$ and $a=3.5$, the desired expectation is $1+7.5/K$, not $8.5$ for every $K$. The replication runs 30,000 independent batches for each $K\in\{1,2,8,64\}$ and records Monte Carlo standard errors. All desired-objective mean checks fall within five reported standard errors. Appendix~\ref{app:trace} gives an unbiased two-probe cross-product correction and explains why unbiased trace estimation does not imply an unchanged training objective after squaring.

A dense Hessian need not be materialized. For diffusion loading with $m$ columns, the exact contraction is a sum of $m$ directional second derivatives; stochastic contraction requires $K$ such probe operations plus loading costs. Network size, parameter differentiation, sample count, optimization iterations, and the probe budget needed for accuracy all enter total cost. There is no model-independent crossover at state dimension 50 or 100.

Finally, the coupled dynamic check has $d\in\{4,8,16\}$, eight dates,
$x_{n+1}=Ax_n+a_n+S\varepsilon_{n+1}$,
$A=.8I+.04\operatorname{Adj}$, $Q=I+.12\mathbf1\mathbf1'/d$, $R=.3I$, and $S=.1I+.02\mathbf1\mathbf1'/d$.
A linear neural actor is improved against a frozen quadratic policy evaluator, and compared with an independent Riccati solution. Off-diagonal transition, payoff, and shock terms are retained. Maximum held-out value errors at the three dimensions are respectively $6.60\times10^{-14}$, $1.55\times10^{-14}$, and $3.32\times10^{-14}$. This is an unconstrained linear--quadratic dynamic regression, not the earlier static quadratic program. Riccati evaluation is faster, and the known quadratic value family makes the test structurally favorable. No constrained sparse-grid superiority or general high-dimensional complexity conclusion follows from it.

\section{Conclusion}
NBO treats policy evaluation, feasible improvement, and independent economic evaluation as different computational objects. The distinction yields explicit value-loss bounds and prevents a critic residual or an actor stationary point from being mistaken for a solution certificate. In the preference--portfolio application, an explicit liquidation contract makes the stochastic model and its numerical transition agree; the re-solved adjustment-cost panel implements a global comparative static for the discounted preference-adjustment budget. Actual neural updates, policy evaluation, diffusion moments, temporal-self recursion, trace estimands, and rival-gradient conventions are now executable.

The economic scope includes recursive utility and strategic interaction without assigning them identical mathematical hypotheses. Further accuracy must be earned through the declared error terms: the current NDU runs certify the finite stochastic economy, while continuous-diffusion refinement remains material. All historical models, derivations, figures, and prior reports remain in the repository, with a retention map distinguishing corrected equations, supplementary material, archival arrays, and executed evidence. This structure permits the next review to assess the mathematics and the numerical economy rather than reconstruct their intended relationship.

\begin{appendix}
\section{Proofs and the scope of the value certificates}
\label{app:proofs}
\subsection{Proof of Proposition~\ref{prop:exact}}
Write $V_j=E(\pi_j)$. The pointwise selector inequality is
$\mathcal H^{I(\pi_j)}[V_j]\geq\mathcal H^{\pi_j}[V_j]$. Comparison and the assumed verification inequality imply $V_{j+1}\geq V_j$. Precompactness of the evaluation image implies boundedness. Therefore at each state-time point the monotone sequence has a limit $\bar V$.

Every subsequence of $V_j$ has a further convergent subsequence in the declared strong norm. Its uniform function limit must be $\bar V$, by monotonicity and pointwise uniqueness of the limit. Consequently the whole sequence converges uniformly to $\bar V$; in fact it converges in the declared strong norm because otherwise a subsequence separated from its unique possible strong cluster point would contradict precompactness. Derivatives of any strong cluster point are the derivatives of its function limit. This observation includes the time derivative, which must be controlled when passing an evaluation equation to the limit.

Let $\pi_{j_m}\to p$ along a policy subsequence. Continuity of evaluation gives $E(p)=\bar V$. Continuity of $I$ at $p$ and of $E$ gives
\[
E(I(p))=\lim_m E(I(\pi_{j_m}))=\lim_m V_{j_m+1}=\bar V.
\]
The policy $I(p)$ attains the Hamiltonian maximum evaluated at $E(p)=\bar V$. Its evaluation at this same value therefore gives
\[
-\bar V_t=\mathcal H^{I(p)}[\bar V]=\sup_{a\in A}\mathcal H^a[\bar V].
\]
The boundary condition passes to the limit in the declared topology. Comparison identifies the value, and verification gives optimality. In particular $E(p)=V^*$, even if $p$ and $I(p)$ differ on ties. No step identifies successive limiting \emph{policies} merely from a small difference of their values.

This proposition concerns an exact operator on a compact, selector-closed policy class. Neither a finite tanh network nor a finite-sample optimizer automatically meets those assumptions. The implemented finite policy is instead assessed by Theorem~\ref{thm:finite}. For recursive drivers the comparison assumptions have to be supplied for that driver and value domain; the proof does not obtain comparison from the word ``recursive.''

\subsection{Proof of Theorem~\ref{thm:continuous}}
Apply discounted It\^o's formula to the test function under any feasible policy $a$ and stop at $\tau$. Taking expectations yields the identity
\begin{equation}
J^a(t,s)-v(t,s)=\mathbb E_{t,s}\left[e^{-\rho(\tau-t)}(g-v)(\tau,S_\tau)
+\int_t^\tau e^{-\rho(r-t)}R_a(v)(r,S_r)\,dr\right].
\label{eq:ito-identity}
\end{equation}
If necessary, localization followed by the integrability hypothesis justifies removal of the localizing sequence. For the enacted policy, $R_\pi(v)\geq-\epsilon$ and the boundary difference is at least $-b$. Hence
$J^\pi\geq v-b-C_\rho(T-t)\epsilon$. For every feasible alternative,
$R_a(v)\leq R_\pi(v)+\eta\leq\epsilon+\eta$ and the boundary difference is at most $b$. Thus
$J^a\leq v+b+C_\rho(T-t)(\epsilon+\eta)$. Taking a supremum over alternatives and subtracting proves \eqref{eq:continuous-bound}. Attainment of that supremum is not required.

The proof does not need differentiability of the true value or differentiability of the policy with respect to its parameters. It does need a legitimate test function, a legitimate stopped process, and bounds valid wherever any competing policy can travel. A residual measured under the actor's occupancy alone need not control an alternative policy. This is why an off-policy uniform audit, or an explicitly proved change-of-measure bound, is essential for the displayed worst-case guarantee.

\paragraph{Boundary and domain qualifications.}
The theorem bounds the stopped payoff as defined, even when the optimal value lacks a classical extension to every boundary face. The NDU diffusion can be degenerate in wealth under $\pi=0$. We do not assert that all boundary points are regular for every controlled diffusion, or that the continuous NDU value belongs to $C^{1,2}$ up to all faces. The scalar payoff at a boundary starting state is settlement by definition; an interior limiting value may require a generalized viscosity boundary interpretation. The theorem's test-function and residual hypotheses would still have to be verified before numerical use. The finite boundary-switched critic is not used to bypass this requirement.

\paragraph{Sampling and covering errors.}
Suppose a residual $R$ is known to be $L_R$-Lipschitz on a compact set and audit nodes form a $d_h$-net. Then
\[
\|R\|_\infty\leq\max_j|R(s_j)|+L_Rd_h.
\]
The same argument applies to a feasible improvement gap if its Lipschitz constant is established. Finite action coverage gives an additional $L_A h_A$ term only when a uniform action Lipschitz bound $L_A$ is known. These inequalities are deterministic covering statements, not consequences of a large sample size. No such global constants have been computed for the NDU neural critic, and no continuum certificate is reported for it.

\subsection{Proof of Theorem~\ref{thm:finite}}
Positivity and row sums at most one imply that $T_n^\pi$ and $T_n$ are monotone and $\delta$-Lipschitz in the sup norm. Let $d_n^\pi=\|\widehat v_n-V_n^\pi\|_\infty$. The evaluation defect gives
$d_n^\pi\leq e_n+\delta d_{n+1}^\pi$, and hence
\[
d_n^\pi\leq\delta^{N-n}b+\sum_{j=n}^{N-1}\delta^{j-n}e_j.
\]
For the optimal value, the residual with respect to $T_n$ is at most $e_n+\eta_n$ in absolute value, because
$\widehat v_n-T_n\widehat v_{n+1}=(\widehat v_n-T_n^\pi\widehat v_{n+1})-(T_n\widehat v_{n+1}-T_n^\pi\widehat v_{n+1})$.
The same backward estimate gives
\[
\|\widehat v_n-V_n^*\|_\infty\leq\delta^{N-n}b+
\sum_{j=n}^{N-1}\delta^{j-n}(e_j+\eta_j).
\]
The triangle inequality proves the result. For exact re-evaluation, both $e$ and $b$ vanish; the code computes each date's maximum feasible gain and accumulates it with the factor $\delta$.

Early liquidation creates no additional continuation coefficient: its payoff, including the fractional-step discount, is already in $r_n^a$. The substochastic surviving row sums remain at most one. Boundary nodes are replaced by their immediate settlement operator in both the optimizer and policy evaluator. The theorem therefore applies to the finite model even though its boundary policy outputs are not economically enacted.

\paragraph{Adding a verified numerical-economy error.}
Suppose, for the same policies and initial states, an independent argument establishes $|J^\pi-J_h^\pi|\leq\zeta$ uniformly and $|V^*-V_h^*|\leq\zeta_*$. Then
$V^*-J^\pi\leq\zeta_*+(V_h^*-J_h^\pi)+\zeta$. This decomposition identifies the missing bridge without assigning it a zero. In particular, a coarse--fine value difference is neither $\zeta$ nor $\zeta_*$ without a further error argument.

\subsection{Proof of Theorem~\ref{thm:cost}}
For each admissible policy, $J_k(\pi)=A(\pi)-kB(\pi)$ is affine and nonincreasing in $k$, since $B\geq0$. The supremum of these affine functions is convex; the pointwise supremum of nonincreasing functions is nonincreasing. At an initial state where optimizers exist, their optimality gives
\begin{align*}
A(\pi_1)-k_1B(\pi_1)&\geq A(\pi_2)-k_1B(\pi_2),\\
A(\pi_2)-k_2B(\pi_2)&\geq A(\pi_1)-k_2B(\pi_1).
\end{align*}
Adding gives $(k_2-k_1)[B(\pi_1)-B(\pi_2)]\geq0$. This holds for any choices of the two optimizers, including multiple optimal policies. Moreover $-B(\pi_k)$ is a subgradient of $V$ at $k$ because
$V(\widetilde k)\geq V(k)-(\widetilde k-k)B(\pi_k)$. At a differentiability point this subgradient equals $V'(k)$. Finally the unrestricted feasible policy set contains the set with $\theta=0$, so its supremum is no smaller.

The argument survives the killed finite approximation because its transition, discounting, and settlement do not depend on $k$, while the per-action cost remains exactly $k$ times its discounted adjustment exposure. It would not apply without modification if changing $k$ also changed the liquidation fee, the feasible action set, the utility normalization, or the law of preference shocks. It says nothing about a pointwise ordering of $\theta$, a portfolio-share ordering, or identification of the model from observed data.

\section{Global action selection and the NDU stochastic operator}
\label{app:ndu}
\subsection{Constrained Hamiltonian maximization}
At a fixed critic jet, the NDU maximand separates across $c,\theta,\pi$. Let their bounds be $(c_-,c_+)$, $(\theta_-,\theta_+)$, and $(\pi_-,\pi_+)$. Consumption maximizes $c^{1-u}/(1-u)-V_X c$, whose second derivative is $-u c^{-u-1}<0$. If $V_X>0$, its optimum is the projection of $V_X^{-1/u}$ onto $[c_-,c_+]$. If $V_X\leq0$, the objective is increasing and the optimum is $c_+$. Adjustment maximizes $V_u\theta-k\theta^2/2$ and therefore selects the projection of $V_u/k$.

Write the portfolio objective as $a_1\pi+a_2\pi^2/2$, where
\[
a_1=.06XV_X+.2(.05)(-.25)XV_{uX},\qquad a_2=.2^2X^2V_{XX}.
\]
If $a_2<0$, select the projected value $-a_1/a_2$. If $a_2>0$, compare the two endpoints. If $a_2=0$, choose the endpoint indicated by the sign of $a_1$, with any feasible action optimal when $a_1=0$. Thus a nonconcave portfolio slice need not create an inaccessible local search problem in this particular model. The finite algorithm instead enumerates its declared action grid; it never labels a stationary neural parameter as a global action optimum.

For a general differentiable concave maximand, concavity gives
$H(a)\leq H(a_0)+\langle\nabla H(a_0),a-a_0\rangle$. Taking a supremum proves \eqref{eq:fwgap}. If an interior maximizer and an $m$-strongly concave objective are known, the familiar gradient-square bound $H^*-H(a_0)\leq\|\nabla H(a_0)\|^2/(2m)$ follows by completing the square. The supporting-hyperplane formulation is preferable at constrained optima. Neither argument transfers concavity from actions to shared network parameters.

\subsection{Local consistency and preference reachability}
Let $b=(\theta,(r+e\pi)X-c)$ and let $C$ denote the covariance matrix of the two-dimensional diffusion. Before the exit correction, \eqref{eq:quadrature} satisfies
\[
\mathbb E\Delta S=b\Delta,\qquad
\operatorname{Cov}(\Delta S)=C\Delta,
\quad C_{uX}=\sigma_u\sigma_X\pi X\rho_{uX}.
\]
The six polynomial tests $1,u,X,u^2,X^2,uX$ give the correct constant and linear generators exactly, and quadratic-generator errors $O(\Delta)$ from the drift outer product. At $(u,X)=(2,1)$ and $(c,\theta,\pi)=(.4,.13,.65)$, the continuous generators are
\[
(0,\ .13,\ -.341,\ .5225,\ -.6651,\ -.553625).
\]
The executable checks step sizes $.1,.05,.025,.0125$. The covariance error is below $10^{-12}$ and the quadratic-generator discrepancy declines at the predicted rate. These tests include, rather than infer, the cross term. At a zero-adjustment interior state, the generator contribution of $u^2$ is $.0025$, not zero.

There is no nearest-neighbor rounding of the next preference state. The mean preference displacement is $\theta\Delta$ and bilinear continuation varies between nodes. A smooth test function linear in $u$ therefore sees every nonzero feasible adjustment even when $|\theta|\Delta$ is less than half a grid spacing. Preference reachability is not inferred from the existence of a control parameter in metadata.

Positive bilinear interpolation on a smooth function has error of order $h_u^2+h_X^2$ under appropriate derivative bounds. Dividing by $\Delta$ produces a generator contribution of order $(h_u^2+h_X^2)/\Delta$. Thus an interior consistency limit requires a coupled choice such as $(h_u^2+h_X^2)/\Delta\to0$, together with vanishing action and time errors. Near the killed boundary, additional exit-time consistency and boundary regularity are needed. The four-branch linear-segment rule is not an exact Brownian-bridge crossing calculation. These qualifications explain why the displayed separate refinements are informative diagnostics but not a continuum error certificate.

\subsection{Boundary cash flows and recorded exits}
For a proposed increment $d$, the exit fraction is the minimum positive fraction at which $s+\lambda d$ reaches a domain face, capped at one. At a crossing the continuation value is $g(t+\lambda\Delta,s+\lambda d)$, the flow receives discount integral
$[1-\exp(-\rho\lambda\Delta)]/\rho$, and there is no continuation policy after settlement. A regression starts at $X=.51$, chooses $c=.8,\pi=0$, and verifies that every branch reaches $X=.5$ with fraction strictly below one and then terminates. The program never changes wealth to $.5$ and permits another consumption date under the pre-exit contract.

The exit ledger distinguishes an exit event from a post-correction state. The adjustment budget is accumulated only until settlement. For a surviving branch interpolated to a boundary node, the next-date boundary settlement is part of the finite kernel. This interpolation effect is a numerical feature, not a physical reflection. The settlement fee is a model primitive; comparing fees 6, 8, and 10 is explicitly labeled a comparison of contracts.

\section{Recursive aggregation, temporal selves, and strategic interaction}
\label{app:recursive}
\subsection{The recursive aggregator and its domain}
For $a=1-1/\psi$, $b=1-\gamma$, and $bV>0$, expand the aggregator as in \eqref{eq:ez}. Its consumption derivative is
\[
f_c(c,V)=\rho c^{a-1}(bV)^{1-a/b}>0.
\]
Its value derivative is
\[
f_V(c,V)=\frac{\rho}{a}(b-a)c^a(bV)^{-a/b}-\frac{\rho b}{a}.
\]
On compact ranges with consumption bounded away from zero and $bV$ bounded away from zero, these derivatives are bounded. This is a local Lipschitz conclusion on the specified range, not a global sign or existence theorem for all values. The transform $V=e^W/b$ enforces $bV>0$, but cannot by itself verify integrability, terminal compatibility, comparison, or transversality.

A policy-wise recursive evaluation solves \eqref{eq:bsde}; in the Markov setting its equation is $v_t+\mathcal L^\pi v+f(c^\pi,v)=0$. When $f$ is Lipschitz in $v$ with constant $L$, standard comparison of a perturbed evaluation with a verified evaluation gives a finite-horizon error envelope of the form
\[
e^{L(T-t)}b+\int_t^T e^{L(r-t)}\epsilon\,dr
\]
under the relevant comparison and integrability hypotheses. A decreasing driver can give a sharper dissipative estimate, but monotonicity must be checked from $f_V$ on the actual range. The time-separable discounted bound in Theorem~\ref{thm:continuous} is not silently applied to an arbitrary Epstein--Zin driver.

\subsection{Executed implicit recursive reference}
The supplementary calculation fixes $\gamma=5$, $\psi=1.5$, and $\rho=.04$ on a deterministic saving problem with ten steps of length $.05$. At each wealth node $x\in[.5,2]$, consumption is $c=mx$ with 80 fractions in $[.01,.6]$, and the next wealth is the stated saving transition in the code. For each action the scalar evaluation equation is
\[
y=I_hV_{n+1}(x')+\Delta f(c,y),\qquad by>0.
\]
The negative-value root is bracketed and solved numerically before maximizing over the action grid. The terminal target is $x^b/b$. Outside the wealth lattice, continuation uses the specified terminal homothetic tail as numerical boundary data at that date. It is not a solution on an unbounded wealth domain. At the central node the recorded value is $-.29658656$ and the selected consumption fraction is the lower bound $.01$. This active bound is reported rather than described as an interior stationary ratio.

The sign-domain and implicit-equation residual tests are actual evaluated quantities. Neural solution error is null because this experiment solves scalar implicit reference equations rather than training a recursive neural policy. The general recursive portfolio model remains part of the paper's operator formulation and analytical derivations; its stationary ratios $.3$ and $.0455$ remain algebraic anchors until their full infinite-horizon verification is supplied.

\subsection{Ordinary continuation and acting-self utility}
For a fixed continuation policy, ordinary log-wealth utility in the two-consumption-date example has form $V_n(w)=A_n\log w+B_n$. Starting from $A_N=1,B_N=0$, the sophisticated self chooses
\[
\alpha_n=\frac{1}{1+\beta\delta A_{n+1}},\quad c_n=\alpha_n w,
\quad A_n=1+\delta A_{n+1},\quad
B_n=\log\alpha_n+\delta[A_{n+1}\log(1-\alpha_n)+B_{n+1}].
\]
Evaluation uses $\delta$, while improvement uses $\beta\delta$. With $\delta=1$ and two decisions, $A_1=2$ and $\alpha_0=1/(1+2\beta)$. Substituting $\beta$ into the evaluation coefficient instead produces the different expression $1/(1+\beta+\beta^2)$. At $\beta=1$ the two formulas coincide, so this special case cannot diagnose the distinction.

A one-shot deviation is evaluated against the same frozen continuation policy whose equilibrium is being claimed. Replacing its continuation with a different grid's optimized value changes the estimand. For a general finite sophisticated-self computation, all feasible actions can be enumerated at each date using the correctly evaluated future policy. For continuous actions a local search gives only a feasible deviation lower bound unless concavity, a global oracle, or an explicit upper-bound calculation closes the search gap.

\subsection{Dynamic games and equilibrium selection}
Let $J_i(\pi_i,\pi_{-i})$ denote player $i$'s payoff with its own objective and transition law. A Nash certificate bounds
\[
\sup_{\widetilde\pi_i}J_i(\widetilde\pi_i,\pi_{-i})-J_i(\pi_i,\pi_{-i}),
\]
not a joint-profit improvement. An actor graph that differentiates through rivals is not the unilateral derivative graph. In the static benchmark $q_i(1-q_i-q_j)$, the derivative with a detached rival is $1-2q_i-q_j$. It vanishes at $q_i=q_j=1/3$, whereas joint optimization selects $1/4$ and leaves a profitable deviation. The executable graph regression verifies both the own gradient and the absence of a rival parameter gradient.

The three-date capacity game has an exogenous Markov demand state and endogenous future capacity. At each state-date pair, enumerate the two players' current feasible output--investment actions using continuation values from the already selected future equilibria. A pure equilibrium is an action pair that is a best response for each player against the other current action and the fixed future policy profile. Every enumerated state has at least one pure equilibrium in this particular game. Lexicographic tie selection is recorded, and the full profile includes asymmetric choices. In a finite-horizon game, this backward best-response construction establishes subgame perfection on the finite state space; no uniqueness follows from it.

For large-population or competitive-limit analysis, the demand and capacity scaling, investment technology, and equilibrium concept must be fixed before taking a limit. Symmetric initialization is neither a proof of equilibrium nor an equilibrium-selection theorem. Those questions are retained as model extensions rather than filled with the historical illustrative $N$-series. No claim about a trained high-dimensional game or a solved competitive limit is attached to the current exact finite game.

\section{Trace estimands and computational cost}
\label{app:trace}
For standard Gaussian $z$ and symmetric $A$, diagonalization and the independent squared-normal moments give
$\mathbb E z'Az=\operatorname{tr}A$ and $\operatorname{Var}(z'Az)=2\|A\|_F^2$.
Independence across $K$ probes yields variance $2\|A\|_F^2/K$ for their mean. Expanding the squared residual then proves \eqref{eq:trace}. For Rademacher probes the diagonal contribution to variance vanishes, so the formula must be changed rather than reused uncritically.

For $K\geq2$, define the cross-product correction
\begin{equation}
U_K=a^2-a\bar q_K+\frac{1}{4K(K-1)}\sum_{i\ne j}q_iq_j.
\label{eq:U}
\end{equation}
Independence of distinct probes gives $\mathbb E U_K=(a-\operatorname{tr}A/2)^2$. The statistic is unbiased but can be negative in a finite batch and may have a different variance and optimization behavior. It is therefore an additional estimand, not a blanket replacement justified by unbiasedness alone. The audit reports its empirical mean and standard error separately. When $A$ depends on network parameters, the uncorrected variance term itself varies with those parameters; a squared stochastic trace can consequently change the expected minimizer.

Let $C_{\rm jet}$ denote the actual cost of a directional second derivative and its parameter gradient for the chosen network. If $\Sigma$ has $m$ columns, exact contraction can be organized as $m$ directional second derivatives, with loading and batching costs. A $K$-probe contraction costs roughly $K C_{\rm jet}$ plus those costs; it is not an unconditional $O(d)$ solver. Total work additionally multiplies by collocation samples, training iterations, and any restarts. Peak memory, warm-up, and accelerator synchronization must be measured for a hardware comparison. The current CPU timers do not measure peak memory, so that metric remains null.

\paragraph{The coupled dynamic regression.}
For the finite-horizon quadratic-cost economy, a policy $a=Wx$ gives a quadratic value $x'Px+c$. Given a frozen next evaluator, the actor minimizes the expected quadratic stage-plus-continuation cost. Gaussian shocks contribute $\operatorname{tr}(PSS')$; the value recursion includes this constant. A matrix-valued linear neural actor is trained against the frozen evaluator by L-BFGS, while an independent Riccati recursion computes the benchmark. Both use identical $A,Q,R,S$, horizon and terminal cost. Correlation is a state-transition covariance here, not a metadata label on a static objective.

The known quadratic family makes this a favorable regression. At dimensions 4, 8, and 16, Riccati times are approximately $.00060,.00040,.00043$ seconds and actor times are $1.319,.0323,.0384$ seconds. The first timing includes startup effects. These numbers neither establish a dimension scaling law nor compare against adaptive sparse grids. The original static coupled quadratic program is retained in the R3 archive as a small linear-algebra check, separate from this dynamic experiment.

\section{Approximation, counterexamples, and optimization limits}
\label{app:limits}
\subsection{A composite objective can bias evaluation}
Consider a one-state terminal evaluation problem whose true value is zero and whose critic is a scalar $v$. The exact evaluation residual is $v$. A joint surrogate $v^2-\lambda v$ has minimizer $v=\lambda/2$, not the evaluated value, for every $\lambda\ne0$. The example isolates the logical issue: optimizing performance and evaluation over the same value parameter need not preserve the evaluation equation. Adding a positive residual weight changes the bias but does not prove exact evaluation. The separated actor has no gradient into that value parameter, so it does not create this particular bias.

Similarly, zero-centering a residual is not a stability theorem. The derivative of $R^2/2$ is $R\nabla R$ and its Jacobian, conditioning, sampling, and step sizes govern a finite optimizer. Value dependence in a recursive driver remains inside this derivative. A zero target does not make the objective independent of the value level.

\subsection{A small squared residual does not select a viscosity solution}
On $(-1,1)$ with boundary values zero, consider $|v'|=1$. The nonnegative distance function $1-|x|$ is the viscosity solution. The function $|x|-1$ has the same boundary values and satisfies $|v'|=1$ almost everywhere, but it fails the viscosity supersolution test at zero: the constant test function $-1$ touches it from below and has $|\phi'|-1=-1<0$.

Smooth approximations
$v_\varepsilon(x)=\sqrt{x^2+\varepsilon^2}-\sqrt{1+\varepsilon^2}$
have exactly the boundary data and converge uniformly to $|x|-1$. Their squared residual $\int_{-1}^1(|v_\varepsilon'|-1)^2dx$ tends to zero by dominated convergence, although the limit is not the viscosity solution. Thus boundary agreement, smooth networks, and an $L^2$ residual do not establish viscosity selection. Monotone, stable, consistent schemes with appropriate comparison can provide a different route \citep{barles1991}; positivity of an interior interpolation alone does not verify all of those hypotheses at a killed boundary.

Smooth nonpolynomial neural networks can approximate a sufficiently smooth function and derivatives on compact subsets under the relevant approximation hypotheses. This does not create derivatives at a kink, ensure global boundary regularity, close a finite policy class under a maximizing selector, or prove that stochastic optimization finds the approximant. The retained derivative-approximation motivation is used only in that conditional sense.

\subsection{Two-timescale statements and suboptimal attractors}
For projected stochastic approximation, a conventional ODE analysis requires bounded iterates, suitable continuity and stability, martingale-difference noise with bounded conditional second moments, step sizes satisfying infinite sums, finite squared sums, and a vanishing slow-to-fast step-size ratio, plus a globally attracting critic equilibrium for each frozen actor. Under the appropriate additional technical conditions, the actor limit set is governed by the corresponding limiting ODE and its internally chain-transitive sets.

This is not a global optimality statement. For $H(a)=-(a^2-1)^2+.2a$, the negative local maximizer is an asymptotically stable isolated stationary point with strictly positive global improvement gap. The audit solves the cubic stationarity equation and evaluates the gap $.3998745872$. Exact critic evaluation and zero noise do not remove that attractor. To conclude optimality one must establish, for the actual actor dynamics, that every relevant limiting set has zero feasible improvement gap, or use an applicable global concavity or gradient-domination argument. The finite runs use Adam and L-BFGS and are assessed by their measured errors; they are not examples of a diminishing-step asymptotic theorem.

\section{Implementation, evidence classes, and reproducibility}
\label{app:repro}
\subsection{Executed parameterizations and data flow}
The continuous Merton critic has one learned log-amplitude and its actor has two learned logits. It uses float64 automatic differentiation through the value with respect to wealth, including the Hessian; critic and actor optimizers are separate. Its independently evaluated policy has positive transversality decay $D$. Analytical controls are used only after training to measure error.

The NDU run is different: time-specific nonlinear networks represent a finite Bellman operator. Its actor maps two normalized states to 125 logits through two 48-unit tanh layers. Its critic maps six features, including log distances to boundary faces, through two 48-unit tanh layers to a scalar adjustment to terminal utility. Next-date critics are frozen, the actor target uses only their continuation predictions, and the current critic is fitted to the enacted actor's policy-evaluation target. The independent NumPy optimizer and evaluator are not imported as training labels.

Training residuals are not out-of-sample certificates because all interior finite states are collocation points. The independent evaluator changes the computational path, not the economic grid. The source contains separate NumPy and PyTorch implementations of the transition and interpolation, whose equality is tested on deterministic data. Cross-grid reference calculations are discretization diagnostics; they have not been relabeled a held-out neural accuracy result.

\subsection{Interpreting metrics}
A value error against the finite optimal reference, a critic error against the evaluated policy, a coarse--fine value difference, and a unilateral gain are different estimands. Policy component differences are retained for consumption, adjustment, and portfolio share. Boundary-node actor outputs are not enacted after settlement; full-grid maxima and interior summaries must be distinguished when interpreting them. A finite exhaustive best-response maximum is a finite-model certificate. A feasible search over only some continuous deviations is merely a lower bound on exploitability.

Execution status and tolerance status are separate fields. A program that completes can miss its target; a timed-out audit has no final solution certificate. A null numerical error has an explicit reason, such as an unmeasured continuous-diffusion error or a recursive reference that did not train a neural policy. A numerical zero is used only for a quantity actually calculated, such as the finite game deviation maximum or a tested monotonicity violation.

The initial failed neural source and its weights, metrics, and raw arrays are retained in the local review bundle. The successful three-seed sources and weights are separate. The attempted twelve-date runs are marked incomplete. Wall-clock times are metadata and are excluded from scientific payload hashes. Source hashes identify the exact executed script; the reachable review commit identifies the historical base, not a fictitious commit containing the subsequently created scripts. The final manifest maps every delivered source to its actual file hash and repository path.

\subsection{Retained material and revision lineage}
The new authoritative entry is \texttt{ECTA\_R4.tex}. The original \texttt{ECTA.tex}, the R3-authoritative \texttt{ECTA\_R2.tex}, its standalone supplement, embedded arrays, earlier reports, and R3 numerical artifacts are left unchanged. Core model formulations, constrained policy derivations, recursive aggregation, temporal-self and game extensions, trace analysis, and theoretical qualifications appear in the current text and these appendices. Historical unexecuted plots are retained as archival objects, not presented again as learned current outcomes. The retention map identifies their source locations and the corresponding corrected discussion. This preserves the scientific record without repeating incompatible equations as simultaneous current specifications.

\end{appendix}
\begin{thebibliography}{99}
\bibitem[Barles and Souganidis(1991)]{barles1991} Barles, G., and P. E. Souganidis (1991): ``Convergence of Approximation Schemes for Fully Nonlinear Second Order Equations,'' \emph{Asymptotic Analysis}, 4, 271--283.
\bibitem[Brumm and Scheidegger(2017)]{brumm2017} Brumm, J., and S. Scheidegger (2017): ``Using Adaptive Sparse Grids to Solve High-Dimensional Dynamic Models,'' \emph{Econometrica}, 85, 1575--1612.
\bibitem[Cai, Fang, and Zhou(2025)]{cai2025} Cai, W., S. Fang, and T. Zhou (2025): ``SOC-MartNet: A Martingale Neural Network for the Hamilton--Jacobi--Bellman Equation without Explicit inf H in Stochastic Optimal Controls,'' arXiv:2405.03169, version 4, March 15.
\bibitem[Duffie and Epstein(1992)]{duffie1992} Duffie, D., and L. G. Epstein (1992): ``Stochastic Differential Utility,'' \emph{Econometrica}, 60, 353--394.
\bibitem[Epstein and Zin(1989)]{epstein1989} Epstein, L. G., and S. E. Zin (1989): ``Substitution, Risk Aversion, and the Temporal Behavior of Consumption and Asset Returns: A Theoretical Framework,'' \emph{Econometrica}, 57, 937--969.
\bibitem[Han and E(2016)]{han2016} Han, J., and W. E (2016): ``Deep Learning Approximation for Stochastic Control Problems,'' arXiv:1611.07422.
\bibitem[Han, Jentzen, and E(2018)]{han2018} Han, J., A. Jentzen, and W. E (2018): ``Solving High-Dimensional Partial Differential Equations Using Deep Learning,'' \emph{Proceedings of the National Academy of Sciences}, 115, 8505--8510.
\bibitem[Jacka and Mijatovi\'c(2015)]{jacka2015} Jacka, S. D., and A. Mijatovi\'c (2015): ``On the Policy Improvement Algorithm in Continuous Time,'' arXiv:1509.09041.
\bibitem[Judd(1998)]{judd1998} Judd, K. L. (1998): \emph{Numerical Methods in Economics}. MIT Press.
\bibitem[Kim, Kim, Kim, and Cho(2025)]{kim2025} Kim, Y., Y. Kim, M. Kim, and N. Cho (2025): ``Neural Policy Iteration for Stochastic Optimal Control: A Physics-Informed Approach,'' arXiv:2508.01718, version 1, August 3.
\bibitem[Lee and Kim(2024)]{lee2024} Lee, J. Y., and Y. Kim (2024): ``Hamilton-Jacobi Based Policy-Iteration via Deep Operator Learning,'' arXiv:2406.10920.
\bibitem[Merton(1969)]{merton1969} Merton, R. C. (1969): ``Lifetime Portfolio Selection under Uncertainty: The Continuous-Time Case,'' \emph{Review of Economics and Statistics}, 51, 247--257.
\bibitem[Qi(2025)]{qi2025} Qi, Q. (2025): ``Neural Foundations of Endogenous Risk Aversion,'' unpublished manuscript, Peking University; motivating manuscript identified in the retained NBO source.
\end{thebibliography}
\end{document}
